\documentclass[10pt,a4paper]{amsart}
\usepackage[margin=19mm]{geometry}
\usepackage{amsmath,amssymb,mathtools}
\usepackage[scr=rsfs,cal=euler]{mathalfa}
\usepackage{xfrac}
\usepackage[unicode,hidelinks,bookmarks=false]{hyperref}
\newcommand{\ie}{i.e.\ }
\newcommand{\cf}{cf.\ }
\newcommand{\resp}{respectively\ }
\newcommand{\Subsection}{\subsection}
\providecommand{\nobreakdash}{\nobreak\hbox{-}}
\setlength{\emergencystretch}{2em}
\multlinegap0pt
\usepackage{amssymb}
\usepackage{tikz}
\usepackage[shortlabels]{enumitem}
\newtheorem{proposition}{Proposition}
\newcommand{\PP}{\mathbb{P}}\newcommand{\A}{\mathbb{A}}\newcommand{\Z}{\mathbb{Z}}\newcommand{\C}{\mathbb{C}}\newcommand{\N}{\mathbb{N}}\newcommand{\T}{\mathbb{T}}
\title{Enumerative Geometry and Tree-Level Gromov--Witten Invariants}
\author{Reginald Anderson, Carrie Frizzell}
\date{Source: arXiv:2501.03232v3 (CC BY 4.0)}
\begin{document}
\maketitle

\noindent\emph{Source and licence.} Enumerative Geometry and Tree-Level Gromov--Witten Invariants. Version arXiv:2501.03232v3 (CC BY 4.0), distributed by arXiv under the Creative Commons Attribution 4.0 International licence, http://creativecommons.org/licenses/by/4.0/, as recorded in the arXiv licence metadata for this exact version and shown on the abstract page https://arxiv.org/abs/2501.03232v3. Full licence notice: \texttt{licenses/arxiv-2501.03232-CC-BY-4.0.txt}.\par\smallskip

\noindent\textit{Editorial scope.} Counted unit: the article's proposition proposition (source0.tex lines 186--188). The article's complete original proof of it is delivered with it (source0.tex lines 190--203, one \texttt{proof} environment of 14 lines). No mathematics is written, repaired or re-derived: the statement and the proof are the article's own lines, in the article's own notation, and no retained line is substituted or re-flowed.  No further source-local context is needed: every cross-reference of every retained line resolves inside the two delivered spans.  Nothing was omitted from any retained span, so the package carries no omission record.  No retained line contains a citation, so the wrapper carries no bibliography and no bibliography entry.  No retained line contains a figure environment, an image-inclusion command or drawing code, so no image is delivered and the package declares no asset dependency.  Editorial formatting only, in addition: an AMS \texttt{amsart} preamble that restates the article's own declarations and macros used by the retained lines, namely the environments \texttt{proposition} on separate counters, together with the standard packages those lines need.  The retained environments print in this document as Proposition 1 in order of appearance, whereas the article prints the same items with the numbering of its own full document.  The complete native source of the article as distributed in the arXiv e-print for this exact version is bundled unchanged as \texttt{sources/171224/source0.tex}.
\begin{proposition} For $Q$ a smooth quintic threefold in $\PP^4$, we have $c_1(TQ)=0$. 
    
\end{proposition}

\begin{proof} 
To show that $c_1(TQ)=0$, we note that the short exact sequence

\[ 0 \rightarrow TQ \rightarrow T\PP^4\bigg|_Q \rightarrow N_Q \rightarrow 0 \] 

becomes \[ 0 \rightarrow TQ \rightarrow T\PP^4\bigg|_Q \rightarrow \mathcal{O}_{\PP^4}(5) \rightarrow 0 \] using adjunction, since $Q$ is smooth. This implies

\begin{align*}
    c(TQ) &= \frac{(1+h\bigg|_Q)^5}{1+5h\bigg|_Q}\\
    &= \frac{1 + 5h\bigg|_Q + 10 (h\bigg|_Q)^2 + 10 (h\bigg|_Q)^3 }{1 + 5h\bigg|_Q}\\
    &= 1 + 10(h\bigg|_Q)^2 - 40 (h\bigg|_Q)^3
\end{align*}

where $h$ is the hyperplane class of $\PP^4$, so that $c_1(TQ)=0$ as claimed. \end{proof} 

\end{document}
